\chapter{Mất Gốc Tư Duy Số Học}
\chapterintro{Có những nỗi sợ Toán không nằm ở bài mới, mà nằm ở chỗ lỗ hổng cũ chưa lấp nổi}{Mất gốc là câu chuyện dài, nên càng cần nói thật gọn và thật đúng.}

\begin{lucbatbox}{Lục bát: Cũ chưa xong, mới đã tới}
Bài xưa em chưa hiểu tròn\
Bài nay đã tới gõ dồn trước sân\
Lớp trên chồng lớp xuống gần\
Như xây mái mới trên nền quá mỏng\
Đến khi gặp một câu cong\
Em nghe cả gốc bên trong rung rinh
\end{lucbatbox}

\begin{tutuyetbox}{Thất ngôn tứ tuyệt: Hổng từ xa}
Không phải hôm nay mới quên\
Mà từ những lớp trước bên em hở dần\
Một khi gốc rễ phân vân\
Cành trên khó đứng bao lần gió qua
\end{tutuyetbox}

\begin{ghichu}
Mất gốc không đáng cười ở bản thân nó. Điều đáng khịa là nhiều bạn biết mình hổng lâu rồi mà vẫn cố lảng tránh thay vì vá lại từ đầu.
\end{ghichu}

\ngancach

\begin{lucbatbox}{Lục bát: Sợ vì không theo kịp}
Bạn bè hiểu đến bước ba\
Em còn đứng ở cửa nhà bước không\
Muốn đi tiếp cũng nặng lòng\
Vì nền dưới chân còn không yên rồi\
Chậm thôi, vá lại từng nơi\
Còn hơn cứ giả mỉm cười hiểu trơn
\end{lucbatbox}

\begin{tutuyetbox}{Thất ngôn tứ tuyệt: Gốc phải vá lại}
Mất gốc thì phải học dần\
Không ai nhảy vọt qua phần trống sâu\
Đi chậm nhưng bước cho lâu\
Còn hơn đứng mãi phía sau ngại ngùng
\end{tutuyetbox}